\documentclass[11pt]{article}

\usepackage[margin=1in]{geometry}
\usepackage[T1]{fontenc}
\usepackage[utf8]{inputenc}
\usepackage{lmodern}
\usepackage{microtype}
\usepackage{setspace}
\usepackage{hyperref}
\usepackage{graphicx}
\usepackage{booktabs}
\usepackage{array}
\usepackage{longtable}
\usepackage{float}
\usepackage{xcolor}
\usepackage{enumitem}
\usepackage{siunitx}
\usepackage{subcaption}

\hypersetup{
    colorlinks=true,
    linkcolor=blue,
    citecolor=blue,
    urlcolor=blue
}

\setstretch{1.08}

\title{\textbf{Distributed Task Queue System:\\Design, Implementation, and Evaluation}}
\author{Joe Caldwell}
\date{\today}

\begin{document}
\maketitle

\section*{Project Artifacts}
\textbf{Source Code Repository (GitHub):} \url{https://github.com/Mojo4Sho1/CSCE5306_Project2_distributed-task-queue}\\
\textbf{Documentation index:} \texttt{docs/\_INDEX.md}\\
\textbf{Benchmark evidence index:} \texttt{results/loadgen/analysis/starter\_matrix\_2026-02-20/EVIDENCE\_INDEX.md}\\
\textbf{Reproducibility runbook:} \texttt{docs/handoff/STARTER\_MATRIX\_REPRODUCIBILITY.md}

\section{Introduction}
This report presents the design, implementation, and evaluation of a distributed task queue system.
The project studies trade-offs in communication overhead, scalability, and operational complexity across two distributed designs under a shared client-facing API contract and locked fairness controls:
\begin{enumerate}[label=(\alph*)]
    \item a microservice architecture that separates functionality into dedicated services, and
    \item a monolith-per-node architecture where each node colocates most system functionality.
\end{enumerate}
Both designs use gRPC and Protocol Buffers, run with in-memory storage, and target at-least-once execution semantics with best-effort, non-preemptive running cancellation in v1.

\section{System Overview}

\subsection{System Design}
The system accepts client jobs, schedules work, executes tasks, and returns results.
In Design A (microservices), the distributed system includes six \emph{functional} containerized nodes:
\begin{enumerate}
    \item Gateway Service
    \item Job Service
    \item Queue Service
    \item Coordinator Service
    \item Worker Service
    \item Result Service
\end{enumerate}

\noindent
Design B uses six identical monolith nodes where each node hosts the public API service and internal job/queue/coordinator/result services in one process.
For parity, both designs expose the same public API methods: \texttt{SubmitJob}, \texttt{GetJobStatus}, \texttt{GetJobResult}, \texttt{CancelJob}, and \texttt{ListJobs}.

\noindent
The load generator is separate benchmarking tooling and is not counted as a functional node in either design.

\noindent
\textbf{Architecture figure placeholder:}
\begin{figure}[H]
    \centering
    \fbox{\parbox{0.92\linewidth}{
    \centering
    \vspace{1.2cm}
    Insert architecture diagram here (microservices and data flow).
    \vspace{1.2cm}
    }}
    \caption{Microservice-based task queue architecture.}
    \label{fig:architecture}
\end{figure}

\subsection{Communication Model}
The communication model is defined by two protobuf namespaces: public \texttt{taskqueue.v1} and internal \texttt{taskqueue.internal.v1}.
The public service is \texttt{TaskQueuePublicService}; in Design A, internal service boundaries are formalized as \texttt{JobInternalService}, \texttt{QueueInternalService}, \texttt{CoordinatorInternalService}, and \texttt{ResultInternalService}.
gRPC provides explicit request/response schemas, stable method names, and compact binary serialization, which keeps client-visible semantics consistent while allowing different internal topologies across Design A and Design B.

\subsection{Functional Requirements}
Table~\ref{tab:functional-requirements} lists the six functional requirements supported by the system.

\begin{table}[H]
\centering
\caption{Functional requirements of the distributed task queue system.}
\label{tab:functional-requirements}
\begin{tabular}{p{0.12\linewidth} p{0.82\linewidth}}
\toprule
\textbf{ID} & \textbf{Requirement} \\
\midrule
FR-1 & \textbf{SubmitJob}: The system accepts a job request and returns a unique job identifier. \\
FR-2 & \textbf{GetJobStatus}: The system returns current job status (queued, running, done, failed, or canceled). \\
FR-3 & \textbf{GetJobResult}: The system returns job output after completion. \\
FR-4 & \textbf{CancelJob}: The system supports cancellation for pending jobs and best-effort cancellation for running jobs. \\
FR-5 & \textbf{ListJobs}: The system lists recent jobs with filtering and pagination support. \\
FR-6 & \textbf{WorkerHeartbeat/FetchWork}: Workers report liveness and request available work from the system. \\
\bottomrule
\end{tabular}
\end{table}
\noindent
FR-6 is an internal control-plane requirement (coordinator/worker path in Design A) and is not part of the client-facing equivalence surface used for strict A/B semantic parity.

\subsection{Non-Functional Requirements}
In addition to functional behavior, the system satisfies the following non-functional requirements:

\begin{table}[H]
\centering
\caption{Non-functional requirements of the distributed task queue system.}
\label{tab:nonfunctional-requirements}
\begin{tabular}{p{0.12\linewidth} p{0.82\linewidth}}
\toprule
\textbf{ID} & \textbf{Requirement} \\
\midrule
NFR-1 & \textbf{API Contract Consistency}: Both architectures expose equivalent client-facing gRPC APIs and status semantics. \\
NFR-2 & \textbf{Scalability Testability}: The system supports controlled load tests across multiple concurrency levels without code changes. \\
NFR-3 & \textbf{Performance Observability}: The evaluation records throughput (requests/sec and jobs/sec), latency percentiles (p50/p95/p99), and end-to-end completion time. \\
NFR-4 & \textbf{Reproducibility}: A clean environment can reproduce deployment and benchmarking using documented Docker Compose and scripts. \\
NFR-5 & \textbf{Baseline Fault Handling}: The coordinator tracks worker heartbeats and marks workers unavailable after timeout so queued jobs are not silently lost. \\
NFR-6 & \textbf{Capacity Fairness}: Both designs run with the same total worker-slot budget during comparisons. \\
\bottomrule
\end{tabular}
\end{table}


\subsection{Requirement-to-Service Coverage}
Table~\ref{tab:fr-service-map} maps each functional requirement to the primary services that implement it.

\begin{table}[H]
\centering
\caption{Functional requirement coverage by service.}
\label{tab:fr-service-map}
\begin{tabular}{p{0.12\linewidth} p{0.78\linewidth}}
\toprule
\textbf{ID} & \textbf{Primary Services} \\
\midrule
FR-1 & Gateway, Job, Queue \\
FR-2 & Gateway, Job, Result \\
FR-3 & Gateway, Result (canonical readiness resolved via Job status) \\
FR-4 & Gateway, Job, Queue, Coordinator \\
FR-5 & Gateway, Job \\
FR-6 & Coordinator, Worker, Queue \\
\bottomrule
\end{tabular}
\end{table}

\section{Evaluation Methodology}

\subsection{Experimental Setup}
\textbf{Execution host:} All benchmark runs were executed on a single local machine (no multi-host cluster).
The host system is a MacBook Air with an Apple M1 chip (8 CPU cores: 4 performance + 4 efficiency) and \SI{16}{GB} of memory, running macOS \texttt{15.6.1}.
Docker Desktop was used to run Linux containers on Apple Silicon; the Docker Engine reports \texttt{OSType=linux} and \texttt{Architecture=aarch64}.\\

\noindent
\textbf{Container runtime:} Docker Engine \texttt{24.0.6} and Docker Compose \texttt{v2.22.0-desktop.2}.\\

\noindent
\textbf{System size:} Each design was deployed with six functional containerized nodes, and the workload generator ran as a separate client process (not counted as a system node).\\

\noindent
\textbf{Workload driver:} A dedicated load-generator client issued controlled gRPC traffic and wrote per-run artifacts under \texttt{results/loadgen/<scenario\_id>/<run\_id>/}.

\subsection{Experimental Setup}
\textbf{Hardware environment:} local host environment used for the 2026-02-20 starter-matrix executions; this draft intentionally treats hardware/software version capture as limited unless explicitly available in execution logs.\\
\textbf{Number of nodes:} 6 functional containerized nodes per design\\
\textbf{Workload driver:} a separate load-generator client issues controlled gRPC traffic and writes per-run artifacts under \texttt{results/loadgen/<scenario\_id>/<run\_id>/}.\\
\textbf{Starter workload specification (locked baseline):}
\begin{itemize}
    \item load levels: low (\texttt{concurrency=6}, \texttt{request\_rate\_rps=5}), medium (\texttt{24}, \texttt{20}), high (\texttt{48}, \texttt{40});
    \item request-mix profiles: \texttt{submit\_heavy}, \texttt{poll\_heavy}, and \texttt{balanced};
    \item timing windows: \texttt{warmup\_seconds=15}, \texttt{measure\_seconds=60}, \texttt{cooldown\_seconds=15};
    \item primary starter parity policy: \texttt{ListJobs=0\%}.
\end{itemize}

\subsection{Compared Designs}
\begin{itemize}
    \item \textbf{Design A (Microservices):} six dedicated functional services (Gateway, Job, Queue, Coordinator, Worker, Result) with public ingress through Gateway.
    \item \textbf{Design B (Monolith-per-node):} six identical nodes, each colocating public and internal services in-process with per-node in-memory state.
\end{itemize}

\noindent
\textbf{API equivalence and parity scope:}
Both designs expose the same client-facing gRPC methods with equivalent schema-level contracts:
\texttt{SubmitJob}, \texttt{GetJobStatus}, \texttt{GetJobResult}, \texttt{CancelJob}, and \texttt{ListJobs}.
Primary parity conclusions in v1 are restricted to \texttt{SubmitJob}, \texttt{GetJobStatus}, \texttt{GetJobResult}, and \texttt{CancelJob}; \texttt{ListJobs} is reported as secondary due to Design B v1 non-global best-effort scope.


\subsection{Metrics}
\begin{itemize}
    \item RPC throughput: completed successful RPCs per second, by method.
    \item Job throughput: terminal jobs/sec (\texttt{DONE}, \texttt{FAILED}, \texttt{CANCELED}).
    \item Latency percentiles: p50/p95/p99 by method, measured client-side from request send to response receipt.
    \item End-to-end completion time: \texttt{submit\_to\_done\_ms = terminal\_timestamp\_ms - accepted\_at\_ms}.
\end{itemize}

\subsection{Fairness Controls}
The benchmark fairness contract is locked by \texttt{docs/spec/fairness-evaluation.md}. The key controls used for this report are:
\begin{itemize}
    \item \textbf{Capacity parity:} \texttt{TOTAL\_WORKER\_SLOTS=6} for both designs (Design A: 1 worker container x 6 slots; Design B: 6 monolith nodes x 1 slot).
    \item \textbf{Ingress parity:} Design A client load enters Gateway only; Design B uses deterministic routing policy with round-robin for empty \texttt{client\_request\_id} submit traffic and key-based owner routing for non-empty keys.
    \item \textbf{State-coherence routing in Design B:} \texttt{GetJobStatus}, \texttt{GetJobResult}, and \texttt{CancelJob} are routed by \texttt{job\_id}; monolith nodes do not forward job-scoped requests in v1.
    \item \textbf{Measurement parity:} same hardware class, workload matrix, warm-up/cooldown policy, run duration, and client-side measurement implementation across both designs.
    \item \textbf{Run protocol:} fixed environment setup, warm-up, steady-state measurement window, cooldown, and repeated runs with fixed seeds/deadlines.
\end{itemize}
\noindent
Reproducibility commands for a minimal A/B rerun, aggregation rerun, and notebook refresh are documented in \texttt{docs/handoff/STARTER\_MATRIX\_REPRODUCIBILITY.md}.

\section{Performance and Scalability Results}
\subsection{Aggregate Starter-Matrix Signals}
The completed starter matrix contains 27 run IDs per design (54 total runs), as summarized in \texttt{results/loadgen/analysis/starter\_matrix\_2026-02-20/starter\_matrix\_summary.md}.
Across the primary parity methods (\texttt{SubmitJob}, \texttt{GetJobStatus}, \texttt{GetJobResult}, \texttt{CancelJob}), mean throughput is matched between designs under fixed pacing, while p95 latency is lower for Design B in most medium/high-load scenarios.

\begin{table}[H]
\centering
\caption{Headline aggregate comparison from starter-matrix summary artifacts.}
\label{tab:headline-aggregate}
\begin{tabular}{p{0.34\linewidth} p{0.26\linewidth} p{0.26\linewidth}}
\toprule
\textbf{Metric} & \textbf{Design A} & \textbf{Design B} \\
\midrule
Run coverage (starter matrix) & 27 runs & 27 runs \\
Global mean p95 across primary methods & 33.615 ms & 24.808 ms \\
Global mean terminal throughput & 2.470988 jobs/s & 5.343827 jobs/s \\
Terminal throughput coefficient of variation & 0.241133 & 0.021774 \\
\bottomrule
\end{tabular}
\end{table}

\subsection{Plot-Based Interpretation}
Two canonical plots were generated from aggregate CSV outputs:
\begin{itemize}
    \item p95 latency by primary method: \texttt{results/loadgen/analysis/starter\_matrix\_2026-02-20/plots/ab\_p95\_latency\_primary\_methods.png}
    \item terminal throughput by scenario: \texttt{results/loadgen/analysis/starter\_matrix\_2026-02-20/plots/ab\_terminal\_throughput\_by\_scenario.png}
\end{itemize}

The latency plot trend is consistent with the delta table in \texttt{starter\_matrix\_ab\_delta\_primary.csv}: the mean p95 latency delta across primary method-scenario rows is \texttt{-8.807 ms} (\texttt{B - A}), indicating lower tail latency for Design B overall.
At low load, latency differences are mixed (for example, \texttt{s\_low\_submit\_heavy} shows slight A advantage), but as load increases, Design B shows persistent p95 improvements, especially on \texttt{SubmitJob}/\texttt{GetJobStatus}.

The terminal-throughput plot shows the largest divergence in medium/high submit-heavy and high balanced scenarios.
For example, from \texttt{starter\_matrix\_terminal\_agg.csv}: \texttt{high\_submit\_heavy} is \texttt{3.00 jobs/s} (A) vs \texttt{11.105556 jobs/s} (B), and \texttt{high\_balanced} is \texttt{3.866666 jobs/s} (A) vs \texttt{11.494445 jobs/s} (B).
At low-load scenarios, both designs are close (around \texttt{0.86--1.43 jobs/s}), which suggests separation appears primarily when sustained offered load increases.

\subsection{Result-Scope Notes}
All numeric values in this section are traceable to:
\begin{itemize}
    \item \texttt{starter\_matrix\_summary.md}
    \item \texttt{starter\_matrix\_method\_agg.csv}
    \item \texttt{starter\_matrix\_terminal\_agg.csv}
    \item \texttt{starter\_matrix\_ab\_delta\_primary.csv}
\end{itemize}

\section{Analysis of System-Design Trade-offs}
The benchmark data in this repository supports the conclusion that the monolith-per-node baseline (Design B) performed better than the microservice baseline (Design A) for this v1 configuration, particularly under medium/high sustained load.
That interpretation is consistent with your intuition: both designs were run with equal worker-slot capacity (\texttt{TOTAL\_WORKER\_SLOTS=6}), but the microservice topology introduces additional networked service boundaries (Gateway, Job, Queue, Coordinator, Result) that can add queueing and coordination overhead under pressure.

\subsection{Why Design B likely won in this dataset}
\begin{itemize}
    \item \textbf{Fewer cross-process hops per request path:} In Design B, public and internal services are colocated in each monolith process, which likely reduces internal RPC path length and scheduling overhead.
    \item \textbf{Lower tail latency under load:} Aggregate p95 values are lower for Design B across primary methods (global mean delta \texttt{-8.807 ms} for \texttt{B - A}), and this gap grows in higher-load scenarios.
    \item \textbf{Higher and more stable terminal throughput:} Design B's global mean terminal throughput is about \texttt{2.16x} Design A (\texttt{5.343827} vs \texttt{2.470988} jobs/s), with lower relative variability (CV \texttt{0.021774} vs \texttt{0.241133}).
\end{itemize}

\subsection{What this does \emph{not} prove}
\begin{itemize}
    \item It does not prove microservices are universally slower; it shows performance in this specific implementation and environment.
    \item It does not isolate root cause to one service bottleneck without deeper per-service profiling (CPU, queue depth, RPC hop timing, lock contention).
    \item It does not provide unlimited external generalization; external validity is bounded because this is fixed-pacing, environment-specific evidence.
\end{itemize}

\subsection{Operational trade-off perspective}
\begin{itemize}
    \item \textbf{Design A strengths:} cleaner ownership boundaries, easier service-level isolation, clearer mutation authority by domain.
    \item \textbf{Design A costs in this run:} more distributed coordination and potential pipeline saturation points as offered load rises.
    \item \textbf{Design B strengths:} simpler deployment shape and lower coordination overhead per request path in the tested baseline.
    \item \textbf{Design B caveats:} correctness depends on deterministic owner routing in v1; some operations such as \texttt{ListJobs} are explicitly treated as secondary parity scope.
\end{itemize}

Overall, for this project's locked v1 benchmark setup, Design B delivered better latency/throughput behavior while Design A retains architectural modularity advantages that may matter for evolvability and isolation goals beyond raw benchmark throughput.

\section{Lessons Learned from AI Tool Usage}
\subsection{How AI Tools Were Used}
AI assistance was used as an implementation copilot and documentation assistant rather than as an autonomous source of truth.
In code work, it was primarily used to accelerate repetitive engineering tasks such as drafting benchmark runner options and artifact-aggregation utilities (\texttt{scripts/loadgen/run\_benchmark\_scaffold.py}, \texttt{scripts/loadgen/aggregate\_starter\_matrix.py}), plus test scaffolding for routing and loadgen contract checks (\texttt{tests/test\_design\_b\_client\_routing.py}, \texttt{tests/test\_loadgen\_contracts.py}).
In operations work, it was used to standardize command flows for Design A/Design B startup, smoke probes, and minimal reproducibility reruns.
In report work, it was used to convert repository artifacts into traceable prose while keeping references tied to explicit files in \texttt{docs/} and \texttt{results/}.

\subsection{Benefits}
The largest benefit was iteration speed.
AI support reduced time spent on boilerplate-heavy edits (argument parsing, repeated validation command formatting, and documentation synchronization across README/spec/handoff files), which made it easier to keep behavior changes and docs in lockstep.
It also improved design-space exploration velocity by quickly producing candidate structures for fairness controls, routing-policy wording, and benchmark artifact schemas before manual tightening.
For this project workflow, the practical gain was not replacing engineering judgment; it was shortening the edit-review-test cycle for routine but high-volume changes.

\subsection{Limitations and Verification}
AI-generated outputs were treated as draft proposals that required explicit verification.
Three safeguards were used consistently:
\begin{itemize}
    \item \textbf{Contract checks against authority docs/proto:} behavior and wording were reconciled with \texttt{proto/taskqueue\_public.proto}, \texttt{proto/taskqueue\_internal.proto}, and locked specs (especially \texttt{docs/spec/fairness-evaluation.md}).
    \item \textbf{Test-based confirmation:} changes were validated with unit tests and live probes (for example \texttt{tests/test\_loadgen\_contracts.py}, \texttt{tests/test\_design\_b\_client\_routing.py}, and integration smoke scripts in \texttt{tests/integration/}).
    \item \textbf{Artifact traceability and reproducibility checks:} report claims were accepted only when numbers/plots mapped to concrete files under \texttt{results/loadgen/analysis/starter\_matrix\_2026-02-20/} and command flows matched \texttt{docs/handoff/STARTER\_MATRIX\_REPRODUCIBILITY.md}.
\end{itemize}
This process limited hallucinated assumptions and kept AI usage inside a human-audited, test-backed workflow.

\section{Reproducibility Appendix (Concise)}
This report's starter-matrix claims can be reproduced without rerunning the full matrix by following \texttt{docs/handoff/STARTER\_MATRIX\_REPRODUCIBILITY.md}.
The minimal flow is:
\begin{enumerate}
    \item start Design A, run one starter scenario, then stop Design A;
    \item start Design B, run one starter scenario, then stop Design B;
    \item regenerate aggregate tables/plots;
    \item refresh \texttt{notebooks/benchmark\_analysis.ipynb}.
\end{enumerate}

\noindent
\textbf{Minimal command flow (from repo root):}
\begin{verbatim}
docker compose -f docker/docker-compose.design-a.yml up --build -d
conda run -n grpc python scripts/loadgen/run_benchmark_scaffold.py \
  --scenario scripts/loadgen/scenarios/starter_matrix/design_a_s_low_balanced.json \
  --output-dir results/loadgen --live-traffic --precheck-health
docker compose -f docker/docker-compose.design-a.yml down --remove-orphans

docker compose -f docker/docker-compose.design-b.yml up --build -d
conda run -n grpc python scripts/loadgen/run_benchmark_scaffold.py \
  --scenario scripts/loadgen/scenarios/starter_matrix/design_b_s_low_balanced.json \
  --output-dir results/loadgen --live-traffic --precheck-health
docker compose -f docker/docker-compose.design-b.yml down --remove-orphans

conda run -n grpc python scripts/loadgen/aggregate_starter_matrix.py \
  --results-root results/loadgen \
  --output-dir results/loadgen/analysis/starter_matrix_2026-02-20
\end{verbatim}

\noindent
\textbf{Primary artifact locations for report traceability:}
\begin{itemize}
    \item evidence index: \texttt{results/loadgen/analysis/starter\_matrix\_2026-02-20/EVIDENCE\_INDEX.md};
    \item aggregate tables: \texttt{starter\_matrix\_method\_agg.csv}, \texttt{starter\_matrix\_terminal\_agg.csv}, \texttt{starter\_matrix\_ab\_delta\_primary.csv};
    \item plots: \texttt{plots/ab\_p95\_latency\_primary\_methods.png}, \texttt{plots/ab\_terminal\_throughput\_by\_scenario.png};
    \item run log: \texttt{results/loadgen/starter\_matrix\_execution\_2026-02-20.log}.
\end{itemize}

\noindent
Interpretation must retain fairness guardrails from \texttt{docs/spec/fairness-evaluation.md}: equal worker-slot budget, Design B owner-routing rules for job-scoped operations, fixed timing windows, and \texttt{ListJobs=0\%} in primary parity scenarios.

\section{Conclusion}
This project demonstrates practical trade-offs between service decomposition and function colocation in distributed systems.
The evaluation provides empirical evidence on throughput, latency, and scalability behavior under controlled workloads.

\end{document}
